% Lösungen Brüche und Dezimalzahlen · Anteil bestimmen · Prüfungs-Fokus MSA (zusammenbau v0.6)
\einheitenkopf{Einheit 1 · Bruch als Anteil}

\erg{1}{a) ja \quad b) $\frac{3}{5}$ \quad c) $5$ von $6$ Teilen}

1c)

\bruchrechteck{5}{6}

\erg{2}{a) $\frac{3}{9}$ \quad b) $\frac{3}{8}$ \quad c) $\frac{4}{12} = \frac{1}{3}$ \quad d) $\frac{4}{16} = \frac{1}{4}$ \quad e) $\frac{13}{25} = \frac{52}{100} = 52\,\%$ \quad f) $\frac{11}{20} = \frac{55}{100} = 55\,\%$ \quad g) $6$ Kästchen \quad h) $10$ Kästchen \quad i) $8$ Kästchen \quad j) $9$ Kästchen \quad k) $20 : 5 \cdot 4 = 16$ Kästchen \quad l) $18 : 3 \cdot 2 = 12$ Kästchen \quad m) $5$ Kästchen \quad n) $4$ Kästchen \quad o) Beide Diagonalen und die beiden Verbindungen gegenüberliegender Seitenmitten teilen das Quadrat in acht gleiche Dreiecke; eines davon kennzeichnen. \quad p) Zwei zueinander senkrechte Durchmesser teilen den Kreis in vier gleiche Viertel; drei davon kennzeichnen. \quad q) Q, denn $\frac{60}{360} = \frac{1}{6}$ \quad r) S, denn $\frac{72}{360} = \frac{1}{5}$}

2g)

\bruchrechteck{6}{20}


2h)

\bruchrechteck{10}{25}


2i)

\bruchrechteck{8}{32}


2j)

\bruchrechteck{9}{36}


2k)

\bruchrechteck{16}{20}


2l)

\bruchrechteck{12}{18}


2m)

\bruchrechteck{5}{12}


2n)

\bruchrechteck{4}{9}


2p)

\kreissektor{270}{}

